% Calibrar tiro a tiro — la cancha de tejo y las cinco estrategias.
%
% Diagrama propio en TikZ (sin librerías extra: círculos, líneas y nodos).
% Se tiñe con el color de la línea (milcGradeDark = verde lima de Pensamiento
% computacional), con milcMostaza para el bocín y los tiros.
%
% Izquierda: vista superior esquemática de la cancha. Tres tiros numerados que
% se van acercando, cada uno con la lectura que deja y el único ajuste que se
% hace en el siguiente. Derecha: las cinco estrategias, en el orden en que se
% usan frente a un reto.
%
% Las anotaciones de dentro enseñan; la síntesis va en el pie de figura vía
% \guiaDiagrama, no aquí.
\begin{tikzpicture}[font=\sffamily,>=stealth]

  % ── La cancha ────────────────────────────────────────────────────────
  % Línea de lanzamiento
  \draw[milcNegro,line width=1pt] (0,-0.9) -- (0,1.1);
  \node[below,font=\scriptsize\bfseries,milcNegro,align=center,text width=1.8cm]
    at (0,-1.0) {línea de\\lanzamiento};

  % Círculo de arcilla y bocín
  \draw[fill=milcNegro!8,draw=milcNegro!45,line width=0.6pt] (5.4,0.1) circle (1.05);
  \draw[fill=milcMostaza,draw=milcNegro,line width=0.8pt] (5.4,0.1) circle (0.17);
  \node[below,font=\scriptsize\bfseries,milcNegro] at (5.4,-1.2) {bocín};

  % Distancia
  \draw[<->,milcNegro!55,line width=0.5pt] (0,1.45) -- (5.4,1.45);
  \node[above,font=\scriptsize\itshape,milcNegro!70] at (2.7,1.45) {17,5 metros};

  % ── Los tres tiros ───────────────────────────────────────────────────
  % 1 · se pasa largo
  \draw[milcGradeDark!45,line width=0.8pt,dashed] (0,0.1) .. controls (3,1.5) .. (6.7,0.5);
  \draw[fill=milcGradeDark!45,draw=none] (6.7,0.5) circle (0.13);
  \node[font=\scriptsize\bfseries,white] at (6.7,0.5) {1};

  % 2 · se abre a la derecha
  \draw[milcGradeDark!70,line width=0.9pt,dashed] (0,0.1) .. controls (3,-1.1) .. (5.9,-0.55);
  \draw[fill=milcGradeDark!70,draw=none] (5.9,-0.55) circle (0.13);
  \node[font=\scriptsize\bfseries,white] at (5.9,-0.55) {2};

  % 3 · pegado al bocín
  \draw[milcGradeDark,line width=1.1pt] (0,0.1) .. controls (3,0.75) .. (5.62,0.28);
  \draw[fill=milcGradeDark,draw=none] (5.62,0.28) circle (0.13);
  \node[font=\scriptsize\bfseries,white] at (5.62,0.28) {3};

  % ── Lecturas: qué dejó cada tiro y qué se ajusta ─────────────────────
  \node[anchor=west,align=left,font=\scriptsize,milcNegro] at (-0.3,-2.25)
    {\textbf{1.} se fue largo \emph{→ ajusta solo la fuerza}};
  \node[anchor=west,align=left,font=\scriptsize,milcNegro] at (-0.3,-2.85)
    {\textbf{2.} se abrió a la derecha \emph{→ ajusta solo la dirección}};
  \node[anchor=west,align=left,font=\scriptsize,milcNegro] at (-0.3,-3.45)
    {\textbf{3.} quedó pegado \emph{→ repite ese tiro, sin cambiar nada}};

  % ── Las cinco estrategias ────────────────────────────────────────────
  \node[align=left,font=\scriptsize\bfseries,milcGradeDark] at (9.1,1.75) {frente a un reto};
  \foreach \i/\y/\texto in {%
    1/0.85/{leer las reglas dos veces},%
    2/-0.15/{probar con un caso pequeño},%
    3/-1.15/{buscar lo que no cambia},%
    4/-2.15/{descartar lo imposible},%
    5/-3.15/{explicarlo en voz alta}}{
    \node[fill=milcGradeDark,text=white,rounded corners=3pt,inner sep=4pt,
          align=left,font=\scriptsize\bfseries,text width=3.3cm] at (9.1,\y) {\i. \texto};
  }

\end{tikzpicture}
